\noindent\textbf{Extended Data Fig.~X | Reconstruction uncertainty in the association between phenomic diversity and diversification.}
\textbf{a,b}, Distributions of Pearson correlation coefficients (\textbf{a}) and two-sided nominal $p$-values (\textbf{b}) between $N_{\mathrm{eff}}^{+}$ and family-level PyRate net diversification rates.
\textbf{c,d}, Corresponding distributions for BAMM.
Net diversification was calculated as $\lambda-\mu$ within each diversification sample.
Each correlation used 13 complete 20-million-year bins spanning 260--0~Ma.
All 1,000 $N_{\mathrm{eff}}^{+}$ stochastic-map reconstructions were paired with all 2,000 retained PyRate trajectories or all 18,001 retained BAMM trajectories, yielding 2,000,000 and 18,001,000 trajectory pairs, respectively.
PyRate trajectories comprised the final 200 saved samples from each of 10 runs; BAMM trajectories were reconstructed from the existing posterior on the full 765-tip phylogeny.
$N_{\mathrm{eff}}^{+}$ was calculated on the 77-taxon phenomic tree.
Complete trajectories were retained during pairing, preserving their temporal structure.
Histograms show percentages of trajectory pairs; $p$-values are displayed on a logarithmic scale.
Dashed grey lines indicate correlations or tests of the mean trajectories, and dotted orange lines mark $p=0.05$.
Annotations show median correlations and central 95\% reconstruction ranges, or percentages of pairs with $p<0.05$.
These ranges describe variation among reconstructions, rather than sampling confidence intervals.
Trajectory pairs are not independent experimental replicates, and the nominal Pearson tests ($n=13$, 11 degrees of freedom) do not correct for temporal autocorrelation.
